% ====================================================================
%  7.5 · REQUERIMIENTOS NO FUNCIONALES (categorizados · Deriva de RN/RF)
% ====================================================================
\begin{longtable}{|C{1.7cm}|L{2.9cm}|L{6.4cm}|L{3.4cm}|}
\hline
\thd{N\textdegree} & \thd{Categoría} & \thd{Descripción del Requerimiento No Funcional} & \thd{Deriva de} \\
\hline
\endfirsthead
\hline
\thd{N\textdegree} & \thd{Categoría} & \thd{Descripción del Requerimiento No Funcional} & \thd{Deriva de} \\
\hline
\endhead
RNF-001 & Seguridad & Las contraseñas deben tener un mínimo de 8 caracteres, con al menos una mayúscula, una minúscula, un dígito y un carácter especial, y no pueden coincidir con el correo de la cuenta. & RF-001, RN-011 \\
\hline
RNF-002 & Seguridad & Las contraseñas se almacenan mediante una función de hash de un solo sentido con sal única por usuario, nunca en texto plano. & RF-001, RN-011 \\
\hline
RNF-003 & Seguridad & La identidad de la cuenta se obtiene exclusivamente del token firmado (JWT), nunca de parámetros de la petición. & RF-001, RF-004, RN-003 \\
\hline
RNF-004 & Seguridad & La clave utilizada para firmar los tokens JWT se almacena fuera del código fuente mediante variables de entorno o secretos del ambiente. Su rotación no requiere modificar el código y se aplica mediante una actualización segura de la configuración y un reinicio controlado del servicio. & RF-001 \\
\hline
RNF-005 & Seguridad & La autorización se evalúa en cada endpoint, con independencia de lo que exponga la interfaz. & RN-003 \\
\hline
RNF-006 & Seguridad & La revocación de la sesión y la vigencia del rol se verifican en cada petición. & RF-004, RN-005, RN-013 \\
\hline
RNF-007 & Seguridad & El tipo real del archivo se valida en el servidor, sin confiar en la extensión declarada. & RF-013, RN-018 \\
\hline
RNF-008 & Seguridad & Las respuestas de recuperación y el mensaje de error de autenticación son idénticos para toda cuenta (no revelan su existencia). & RF-001, RF-002, RN-010 \\
\hline
RNF-009 & Seguridad & El token del enlace de recuperación se genera con un generador de números aleatorios criptográficamente seguro. & RF-002, RN-010 \\
\hline
RNF-010 & Seguridad & El contenido recuperado del corpus se trata como datos dentro del prompt, nunca como instrucciones (protección ante \emph{prompt injection}). & RF-024, RN-033, RN-034 \\
\hline
RNF-011 & Integridad & La unicidad documental por contenido emplea el algoritmo de hash SHA-256. & RF-014, RN-024 \\
\hline
RNF-012 & Integridad & La transferencia del rol SuperAdmin se ejecuta en una transacción única; un fallo parcial revierte toda la operación. & RF-008, RN-009 \\
\hline
RNF-013 & Integridad & La unicidad del SuperAdmin se garantiza mediante una restricción de integridad en la base de datos, además de la validación en la aplicación. & RN-002 \\
\hline
RNF-014 & Integridad & El modelo de datos distingue entre un valor no determinado y un valor igual a cero. & RN-031, RN-032 \\
\hline
RNF-015 & Integridad & El contenido de los documentos se procesa y almacena en UTF-8, preservando tildes y la letra «ñ» sin alteración. & RF-012, RF-016 \\
\hline
RNF-016 & Rendimiento & El sistema debe procesar la ingesta e indexación de un documento de hasta 50 MB dentro del tiempo máximo validado mediante pruebas de rendimiento, considerando la latencia y los límites del proveedor externo de embeddings. & RF-015 \\
\hline
RNF-017 & Rendimiento & El sistema debe registrar por separado el tiempo de procesamiento interno y el tiempo de respuesta del proveedor externo. Bajo condiciones normales, la respuesta completa del asistente no debe superar el límite definido durante las pruebas de rendimiento. & RF-023 \\
\hline
RNF-018 & Rendimiento & Las llamadas al proveedor de IA tienen un tiempo máximo de espera de 1 minuto, tras el cual la operación se reporta como indisponible. & RF-023 \\
\hline
RNF-019 & Disponibilidad & La indisponibilidad del proveedor de IA no degrada los módulos de gestión, ingesta, reportes de prospección ni auditoría. & RF-024, RN-033 \\
\hline
RNF-020 & Capacidad & El sistema admite y procesa documentos de hasta 50 MB por cada operación de ingesta. & RF-013, RN-021 \\
\hline
RNF-021 & Disponibilidad & El sistema ofrece una disponibilidad de al menos el 95\,\% en horario laboral (lun–vie, 8:00–20:00, hora de Perú). & Transversal \\
\hline
RNF-022 & Trazabilidad & Los registros de auditoría se almacenan en estructuras de solo inserción (\emph{append-only}), sin actualización ni eliminación. & RN-044 \\
\hline
RNF-023 & Trazabilidad & El registro de auditoría conserva sus entradas durante todo el periodo de operación del sistema. & RN-044 \\
\hline
RNF-024 & Trazabilidad & La fecha y hora de la auditoría se almacenan en UTC y se convierten a \emph{America/Lima} solo en la presentación. & RN-043 \\
\hline
RNF-025 & Trazabilidad & Todo fragmento indexado conserva la referencia a su documento de origen y a su ubicación dentro de él. & RF-016, RN-034 \\
\hline
RNF-026 & Mantenibilidad & El catálogo de 40 códigos GRI se carga mediante un script de inicialización versionado, garantizando el mismo catálogo en todos los ambientes. & RN-027 \\
\hline
RNF-027 & Mantenibilidad & El sistema integra al proveedor de IA mediante una interfaz estandarizada, de modo que cambiar de modelo o proveedor no obliga a rediseñar. & RF-023 \\
\hline
RNF-028 & Usabilidad & Durante la ingesta síncrona el sistema mantiene informado al usuario del progreso de la operación. & RF-015, RN-026 \\
\hline
RNF-029 & Usabilidad & Los mensajes de estado y de rechazo se presentan en español claro y accionable, indicando el motivo. & RF-013 \\
\hline
RNF-030 & Continuidad & Se realizan copias de seguridad diarias de la base de datos (incluido el índice vectorial) y de los documentos \texttt{.md} originales. & RF-016 \\
\hline
\end{longtable}
